%======================================================================
\chapter{Enhancements}
\label{ch:improve}
%======================================================================

Seven changes take the simulator from a demonstration to something defensible. The first is
a security fix; the next four replace invented numbers with physics or accounting; the last
two add capability the base paper only describes.

\begin{table}[H]
\centering
\caption{Summary of the seven enhancements.}
\label{tab:improvements}
\small
\begin{tabularx}{\textwidth}{@{}L{7mm} L{34mm} L{40mm} X@{}}
\toprule
& \textbf{Enhancement} & \textbf{Replaces} & \textbf{Why}\\
\midrule
E1 & AES-256-GCM authenticated encryption & Repeating-key XOR &
  Confidentiality \emph{and} integrity; matches the paper's stated use of AES\\
E2 & Thorp absorption model & \code{random.uniform(0.04,0.08)} &
  Delay becomes a function of distance and frequency\\
E3 & 15\,\% Bernoulli loss with 3 retries & Assumed perfect delivery &
  Tests the protocol under the channel's defining property\\
E4 & Per-node battery depletion & Static energy constant &
  Shows which nodes wear out and how fast\\
E5 & Random-waypoint AUV mobility & Static nodes &
  Link distance, hence delay, varies per round as it would in the sea\\
E6 & $z$-score anomaly detection & Nothing &
  Detects delay-injection attacks the paper names but does not counter\\
E7 & Comparison charts and new metrics & 3 basic plots &
  Benchmarks against the five prior schemes; adds throughput\\
\bottomrule
\end{tabularx}
\end{table}

%----------------------------------------------------------------------
\section{E1 --- Replacing XOR with AES-256-GCM}
\label{sec:aes}

This is the most important change in the project, so it gets the most space.

\subsection{Why the XOR cipher had to go}

A repeating-key XOR cipher encrypts by exclusive-or-ing each plaintext byte with a byte of
the key, cycling the key when it runs out. Decryption is the same operation, because
$x \oplus k \oplus k = x$. It has two fatal properties.

\paragraph{Attack 1: known plaintext recovers the key immediately.}
The attacker knows the message format --- \code{"<id>|<nonce>|<timestamp>"} --- and the
identity field is a public SHA-256 hex string she can read off the network during
registration. Given any plaintext byte $p_i$ and its ciphertext byte $c_i$:
\begin{equation}
k_{i \bmod |k|} = p_i \oplus c_i
\end{equation}
Recovering the whole key needs only as many known bytes as the key is long. She then
decrypts every past and future message on that link.

\paragraph{Attack 2: key reuse leaks the relationship between messages.}
The simulator derives one ECDH key per node pair and reuses it for every message on that
link. For two messages under the same key:
\begin{equation}
C_1 \oplus C_2 = (P_1 \oplus K) \oplus (P_2 \oplus K) = P_1 \oplus P_2
\end{equation}
The key vanishes. Because the two plaintexts share a rigid format --- the same 64-character
identity prefix, the same separators, timestamps differing only in their final digits ---
the XOR of the two plaintexts is highly structured, and standard crib-dragging recovers
both.

\paragraph{And there is no integrity at all.}
XOR is malleable. Flipping bit $j$ of the ciphertext flips bit $j$ of the plaintext. The
attacker cannot read the message, but she can \emph{change} it in a controlled way, and the
receiver has no means of noticing. She can alter the nonce, or move the timestamp, without
knowing the key.

\begin{figure}[H]
\centering
\resizebox{\textwidth}{!}{%
\begin{tikzpicture}[x=1cm,y=1cm,font=\scriptsize]
%======================= LEFT: XOR ==============================
\node[font=\small\bfseries\sffamily,text=coral,anchor=west] at (0,0.9) {BEFORE \textemdash\ repeating-key XOR};
\draw[draw=coral,rounded corners=2pt] (-0.25,-7.7) rectangle (7.15,0.55);
\node[box,anchor=west,align=left,minimum width=64mm,fill=white] (x1) at (0,0)
 {plaintext \ \code{"3147de24...|48231|1772451003.4"}};
\node[box,anchor=west,align=left,minimum width=64mm,fill=white] (x2) at (0,-1.2)
 {key \ \ \ \ \ \ \code{"9d2f7c0a..."} \ \ \textcolor{coral}{repeated to length}};
\node[box,anchor=west,align=left,minimum width=64mm,fill=mist] (x3) at (0,-2.4)
 {$c_i = p_i \oplus k_{i \bmod |k|}$ \ \ for every byte};
\node[box,anchor=west,align=left,minimum width=64mm,fill=white] (x4) at (0,-3.6)
 {ciphertext \ \ same length as plaintext, no tag};
\foreach \a/\b in {x1/x2,x2/x3,x3/x4} \draw[flow,draw=coral] (\a.south) -- (\b.north);
\node[box,anchor=west,align=left,minimum width=64mm,fill=palecoral,draw=coral] (x5) at (0,-5.3)
 {\textbf{\textcolor{coral}{Attack A \textemdash\ known plaintext}}\\
  attacker knows the ID prefix (it is public)\\
  $k_i = p_i \oplus c_i$ \ \ $\Rightarrow$ \ \textbf{whole key recovered}};
\node[box,anchor=west,align=left,minimum width=64mm,fill=palecoral,draw=coral] (x6) at (0,-6.9)
 {\textbf{\textcolor{coral}{Attack B \textemdash\ bit flipping}}\\
  flip ciphertext bit $j$ $\Rightarrow$ plaintext bit $j$ flips\\
  \textbf{receiver cannot tell}};
\draw[flow,draw=coral] (x4.south) -- (x5.north);
%======================= RIGHT: AES-GCM =========================
\node[font=\small\bfseries\sffamily,text=kelp,anchor=west] at (8.0,0.9) {AFTER \textemdash\ AES-256-GCM};
\draw[draw=kelp,rounded corners=2pt] (7.75,-7.7) rectangle (15.15,0.55);
\node[box,anchor=west,align=left,minimum width=64mm,fill=white] (g1) at (8.0,0)
 {plaintext \ \code{"3147de24...|48231|1772451003.4"}};
\node[box,anchor=west,align=left,minimum width=64mm,fill=white] (g2) at (8.0,-1.2)
 {key \ \ \ \ \ \ 256 bits from ECDH \ \ \ \ IV \ \ \textcolor{kelp}{12 fresh random bytes}};
\node[box,anchor=west,align=left,minimum width=64mm,fill=mist] (g3) at (8.0,-2.4)
 {CTR mode: $c_i = p_i \oplus \mathrm{AES}_K(\mathrm{IV}\,\|\,\mathrm{counter}_i)$\\
  \textcolor{kelp}{keystream differs per message because the IV does}};
\node[box,anchor=west,align=left,minimum width=64mm,fill=mist] (g4) at (8.0,-3.8)
 {GHASH over the ciphertext $\Rightarrow$ 128-bit tag $\tau$};
\node[box,anchor=west,align=left,minimum width=64mm,fill=paleamber,draw=amber] (g5) at (8.0,-5.0)
 {wire: \ \code{[ IV 12 B ][ ciphertext ][ tag 16 B ]}};
\foreach \a/\b in {g1/g2,g2/g3,g3/g4,g4/g5} \draw[flow,draw=kelp] (\a.south) -- (\b.north);
\node[box,anchor=west,align=left,minimum width=64mm,fill=palekelp,draw=kelp] (g6) at (8.0,-6.9)
 {\textbf{\textcolor{kelp}{Attack A defeated}} \ known plaintext yields one keystream block,\\
  useless for the next message (new IV)\\
  \textbf{\textcolor{kelp}{Attack B defeated}} \ any flipped bit breaks $\tau$; decrypt raises};
\draw[flow,draw=kelp] (g5.south) -- (g6.north);
\end{tikzpicture}}
\caption[XOR versus AES-GCM]{\textbf{The cipher replacement, and what each attack does to
each scheme.} The structural difference is on the right: a fresh random IV per message means
the keystream never repeats even though the ECDH key is long-lived, and the GHASH tag makes
tampering detectable rather than silent. Both are absent from XOR by construction.}
\label{fig:xor-vs-aes}
\end{figure}

\subsection{What AES-GCM provides}

AES-GCM \cite{gcm} is an \emph{authenticated encryption with associated data} (AEAD)
construction. It does two jobs in one pass:

\begin{itemize}
\item \textbf{Confidentiality} via AES in Counter (CTR) mode. AES-256 encrypts a counter
      block, and the resulting keystream is XOR-ed with the plaintext. Note that the final
      combining step \emph{is} an XOR --- the difference from our old code is that the
      keystream is the output of a block cipher and never repeats, rather than a short key
      cycled over and over.
\item \textbf{Integrity} via GHASH, a polynomial hash over the ciphertext in
      $\mathrm{GF}(2^{128})$, producing a 128-bit authentication tag. The receiver
      recomputes the tag and compares. Any modification --- one bit --- makes the tags
      differ, and decryption fails before any plaintext is released.
\end{itemize}

The 96-bit random IV is what breaks the key-reuse problem. The same long-lived ECDH key can
protect any number of messages, because each message's keystream is derived from a different
IV.

\begin{figure}[H]
\centering
\begin{tikzpicture}[x=1cm,y=1cm,font=\scriptsize]
% wire format
\fill[paleamber] (0,0) rectangle (2.6,0.85);
\draw[draw=amber,thick] (0,0) rectangle (2.6,0.85);
\node at (1.3,0.42) {\textbf{IV} \ 12 bytes};
\fill[mist] (2.6,0) rectangle (10.4,0.85);
\draw[draw=navy,thick] (2.6,0) rectangle (10.4,0.85);
\node at (6.5,0.42) {\textbf{ciphertext} \ \ same length as the plaintext (55 bytes here)};
\fill[palekelp] (10.4,0) rectangle (13.4,0.85);
\draw[draw=kelp,thick] (10.4,0) rectangle (13.4,0.85);
\node at (11.9,0.42) {\textbf{tag} \ 16 bytes};
% brace
\draw[decorate,decoration={brace,amplitude=5pt,mirror},draw=slate] (0,-0.15) -- (13.4,-0.15);
\node[font=\scriptsize,text=slate] at (6.7,-0.75) {83 bytes $=$ 664 bits transmitted};
% annotations
\node[font=\tiny,text=amber,align=center,anchor=north] at (1.3,-1.15) {fresh per message;\\ sent in clear\\ (it is not secret)};
\node[font=\tiny,text=navy,align=center,anchor=north] at (6.5,-1.15) {AES-256-CTR keystream $\oplus$ plaintext};
\node[font=\tiny,text=kelp,align=center,anchor=north] at (11.9,-1.15) {GHASH over ciphertext;\\ checked \emph{before}\\ decryption is trusted};
\end{tikzpicture}
\caption[AES-GCM wire format]{\textbf{What actually goes on the wire.} The IV is prepended
so the receiver can reconstruct the keystream; it is public by design. The 28 bytes of
overhead (IV plus tag) are the price of authenticated encryption, and at 10\,kbps cost
22\,ms per message --- an acceptable trade for detecting tampering.}
\label{fig:gcm-wire}
\end{figure}

\subsection{The code change}

\begin{lstlisting}[style=bad,caption={Before: the broken cipher.}]
def encrypt(message, key):
    return ''.join(chr(ord(c) ^ ord(key[i % len(key)])) for i, c in enumerate(message))

def decrypt(cipher, key):
    return ''.join(chr(ord(c) ^ ord(key[i % len(key)])) for i, c in enumerate(cipher))
\end{lstlisting}

\begin{lstlisting}[style=py,caption={After: AES-256-GCM (\code{uwc\_simulation.py},
lines 50--58).}]
from cryptography.hazmat.primitives.ciphers.aead import AESGCM

def aes_gcm_encrypt(plaintext, key_hex):
    key   = bytes.fromhex(key_hex[:64])   # 256-bit key from the ECDH digest
    nonce = os.urandom(12)                # 96-bit random IV, fresh per message
    ct    = AESGCM(key).encrypt(nonce, plaintext.encode(), None)
    return nonce + ct                     # prepend the IV for the receiver

def aes_gcm_decrypt(ciphertext, key_hex):
    key = bytes.fromhex(key_hex[:64])
    return AESGCM(key).decrypt(ciphertext[:12], ciphertext[12:], None).decode()
\end{lstlisting}

\code{AESGCM.decrypt} raises \code{InvalidTag} if the tag does not verify. The caller in
\code{authenticate()} catches it and fails the hop:

\begin{lstlisting}[style=good,caption={Tamper detection at the call site (lines 188--192).}]
try:
    dec = aes_gcm_decrypt(ct, rk)
except Exception:
    print(f"  {r}: Decryption FAILED (tampered message)")
    return False
\end{lstlisting}

\begin{okbox}
\textbf{Relationship to the base paper.} The paper writes each message as
$E\{\PU_{recv}, \langle \cdot \rangle\}$ --- encryption under the recipient's public key ---
and separately states in its energy analysis that the symmetric primitive is AES-128, at
3.2\,\uJ{} per operation. Our construction is the standard hybrid realisation of that:
ECDH under the recipient's public key derives the session key, and AES-256-GCM protects the
payload. This is what public-key encryption schemes such as ECIES do internally. We use
AES-256 rather than AES-128 because the ECDH digest is already 256 bits, and on this
hardware the cost difference is negligible relative to the acoustic transmission.
\end{okbox}

\begin{table}[H]
\centering
\caption{Property-by-property comparison of the two ciphers.}
\label{tab:cipher-cmp}
\small
\begin{tabularx}{\textwidth}{@{}L{44mm} L{40mm} X@{}}
\toprule
\textbf{Property} & \textbf{Repeating-key XOR} & \textbf{AES-256-GCM}\\
\midrule
Confidentiality      & Broken by known plaintext & 256-bit security\\
Integrity            & None                      & 128-bit GHASH tag\\
Tamper detection     & Impossible                & \code{InvalidTag} on any change\\
Key reuse across messages & Catastrophic ($C_1\oplus C_2 = P_1\oplus P_2$) & Safe, given a fresh IV\\
Malleability         & Fully malleable (bit flipping) & Non-malleable\\
Ciphertext expansion & 0 bytes                   & 28 bytes (12 IV $+$ 16 tag)\\
Standardised         & No                        & NIST SP 800-38D\\
Energy per operation & negligible                & $\approx 3.2\,\mu$J (paper's figure)\\
\bottomrule
\end{tabularx}
\end{table}

%----------------------------------------------------------------------
\section{E2 --- The Thorp acoustic absorption model}
\label{sec:thorp}

Sound in seawater does not merely spread out; it is \emph{absorbed}, and the absorption
depends strongly on frequency. Thorp's empirical model \cite{thorp} gives the absorption
coefficient in decibels per kilometre for a frequency $f$ in kHz:

\begin{equation}
\alpha(f) = \frac{0.11 f^{2}}{1+f^{2}} + \frac{44 f^{2}}{4100+f^{2}}
          + 2.75\times10^{-4} f^{2} + 0.003
\label{eq:thorp}
\end{equation}

The four terms correspond to boric acid relaxation, magnesium sulphate relaxation, pure
water viscosity, and a low-frequency residual. Typical underwater modems operate between
10 and 25\,kHz; our simulation uses 25\,kHz.

The one-way delay of a hop is then propagation plus an absorption-derived jitter term plus
a multipath spread term:

\begin{lstlisting}[style=py,caption={Physical delay model (lines 66--85).}]
SOUND_MPS = 1500.0   # speed of sound in seawater

DISTS = {   # realistic inter-node distances, metres
    ("U1","S1"):150,  ("U2","S1"):200,  ("S1","B1"):800,  ("S1","B2"):850,
    ("B1","SAT1"):1000, ("B2","SAT2"):1000, ("SAT1","BS"):500, ("SAT2","BS"):500,
}

def thorp_db_per_km(f_khz):
    f2 = f_khz ** 2
    return 0.11*f2/(1+f2) + 44*f2/(4100+f2) + 2.75e-4*f2 + 0.003

def acoustic_delay(s, r, freq_khz=25.0, mob_off=0.0):
    key  = (s,r) if (s,r) in DISTS else (r,s)
    dist = max(10, DISTS.get(key, 500) + abs(mob_off))   # mobility changes the distance
    prop = dist / SOUND_MPS                              # geometric propagation
    jit  = thorp_db_per_km(freq_khz) * (dist/1000) * 1e-4
    mp   = abs(random.gauss(0, 0.005))                   # multipath spread
    return prop + jit + mp
\end{lstlisting}

\begin{table}[H]
\centering
\caption{Per-hop distances and resulting propagation delays.}
\label{tab:delays}
\small
\begin{tabular}{@{}l r r@{}}
\toprule
\textbf{Hop} & \textbf{Distance} & \textbf{Propagation delay}\\
\midrule
UWS $\to$ SUB  & 150--200\,m & 0.10--0.13\,s\\
SUB $\to$ BUOY & 800--850\,m & 0.53--0.57\,s\\
BUOY $\to$ SAT & 1000\,m     & 0.67\,s\\
SAT $\to$ BS   & 500\,m      & 0.33\,s\\
\midrule
\textbf{Full chain} & \textbf{2450--2550\,m} & \textbf{$\approx$1.63--1.70\,s one way}\\
\bottomrule
\end{tabular}
\end{table}

Note the consequence: the end-to-end round trip is over three seconds. This is precisely why
the protocol authenticates per hop --- a $\Delta T$ of 5\,s is comfortable for a single hop
but would be marginal end to end.

%----------------------------------------------------------------------
\section{E3 --- Packet loss with retransmission}

Stojanovic \cite{stojanovic} measured 10--30\,\% packet loss in open-water acoustic
channels. We model this as an independent Bernoulli trial per transmission at a
conservative 15\,\%, with up to three attempts.

\begin{lstlisting}[style=py,caption={Lossy channel wrapper (lines 92--104).}]
LOSS_RATE = 0.15
MAX_RETRY = 3

def send_with_loss(fn, s, r, ls):
    for attempt in range(1, MAX_RETRY+1):
        if random.random() < LOSS_RATE:
            ls["lost"] += 1
            print(f"  [LOSS] {s}->{r} attempt {attempt}/{MAX_RETRY}")
            continue
        return fn(s, r)
    ls["failed"] += 1
    print(f"  [FAIL] {s}->{r} link down after {MAX_RETRY} retransmissions")
    return False
\end{lstlisting}

\begin{figure}[H]
\centering
\begin{tikzpicture}[x=1cm,y=1cm,font=\scriptsize]
\node[node,minimum width=14mm] (s) at (0,0) {S1};
\node[node,minimum width=14mm] (r) at (10.4,0) {B2};
% attempt 1
\draw[->,draw=coral,thick,dashed] (0.75,0.75) -- (5.2,0.75);
\node[font=\tiny,text=coral] at (6.9,0.75) {attempt 1 \textemdash\ LOST \ ($p=0.15$)};
\node[font=\tiny,text=coral] at (5.55,0.75) {$\times$};
% attempt 2
\draw[->,draw=coral,thick,dashed] (0.75,0.15) -- (5.2,0.15);
\node[font=\tiny,text=coral] at (6.9,0.15) {attempt 2 \textemdash\ LOST \ ($p=0.15$)};
\node[font=\tiny,text=coral] at (5.55,0.15) {$\times$};
% attempt 3
\draw[->,draw=kelp,line width=1.1pt] (0.75,-0.45) -- (9.65,-0.45);
\node[font=\tiny,text=kelp,fill=white,inner sep=1pt] at (5.2,-0.25) {attempt 3 \textemdash\ DELIVERED};
% probabilities
\node[box,anchor=west,align=left,fill=mist,minimum width=58mm] at (-0.2,-1.7)
 {Per-hop success after 3 attempts: \ $1-0.15^{3} = 99.66\,\%$\\
  Full 4-hop chain: \ $(1-0.15^{3})^{4} = 98.65\,\%$\\
  Expected transmissions per hop: \ $\approx 1.18$};
\node[box,anchor=west,align=left,fill=palecoral,draw=coral,minimum width=52mm] at (6.4,-1.7)
 {If all 3 attempts fail the hop is declared\\ down and \code{auth\_path} returns
  \textsc{false};\\ the round is recorded as a failure.};
\end{tikzpicture}
\caption[Retransmission on a lossy hop]{\textbf{Three attempts per hop.} With independent
15\,\% loss, three attempts raise per-hop reliability to 99.66\,\% and end-to-end
reliability across all four hops to 98.65\,\%. This is what makes the protocol usable on a
channel that drops one packet in seven.}
\label{fig:loss}
\end{figure}

%----------------------------------------------------------------------
\section{E4 --- Per-node battery depletion}

Each node begins with a 1000\,mAh cell at 3.3\,V:
\begin{equation}
E_0 = 1000\ \mathrm{mAh} \times 3.3\ \mathrm{V} = 3.3\ \mathrm{J} = 3\,300\,000\ \mu\mathrm{J}
\end{equation}

Every operation now debits the ledger, rather than being counted in the abstract.

\begin{lstlisting}[style=py,caption={Energy accounting (lines 135--145).}]
BAT_uJ = 3_300_000.0
AUTH_uJ, TX_uJ, RX_uJ = 48.8, 50.0, 36.0
batt = {n: BAT_uJ for g in nodes for n in nodes[g]}

def use_energy(node, op="auth"):
    cost = {"auth": AUTH_uJ, "tx": TX_uJ, "rx": RX_uJ}.get(op, AUTH_uJ)
    batt[node] -= cost
    if batt[node] <= 0:
        print(f"  [DEAD] {node} battery depleted!")
        return False
    return True
\end{lstlisting}

The 48.8\,\uJ{} authentication figure is Equation~\ref{eq:energy} from the base paper. The
transmit and receive figures reflect acoustic modem power, and they dominate: a single hop
costs 50 (TX) $+$ 36 (RX) $+$ 48.8 (auth) $= 134.8$\,\uJ, of which the cryptography is only
36\,\%. This is worth stating plainly, because it reframes the optimisation target ---
\emph{reducing message size matters more than reducing cipher cost}, which is exactly why
the paper's 2112-bit figure is its most valuable result.

%----------------------------------------------------------------------
\section{E5 --- AUV mobility}

Submarines and AUVs move at roughly 2\,m/s. Between authentication rounds we drift the
mobile nodes (U1, U2, S1) by up to $\pm10$\,m, clamped to $\pm50$\,m from the nominal
position, following the random-waypoint family surveyed by Camp \emph{et al.}
\cite{camp}.

\begin{lstlisting}[style=py,caption={Mobility update (lines 163--168).}]
mob_off = {n: 0.0 for g in nodes for n in nodes[g]}

def upd_mob():
    for n in ["U1", "U2", "S1"]:
        mob_off[n] += random.uniform(-10, 10)
        mob_off[n]  = max(-50, min(50, mob_off[n]))
\end{lstlisting}

The offset feeds straight into \code{acoustic\_delay}, so link delay varies from round to
round exactly as it would at sea. This is the source of the non-monotonic shape in the delay
plot of Chapter~\ref{ch:results}, and it is a feature, not noise.

%----------------------------------------------------------------------
\section{E6 --- Anomaly detection on the delay stream}

The base paper names \emph{delay injection} in its threat model: an adversary deliberately
slows packets to push a legitimate exchange outside its $\Delta T$ window and force a
desynchronisation. It proposes an adaptive $\Delta T$ as the countermeasure but does not
implement a detector. We added one.

Every observed per-hop delay is recorded. A delay more than 2.5 standard deviations from the
mean is flagged:

\begin{equation}
z_i = \frac{d_i - \mu}{\sigma}, \qquad \text{flag if } |z_i| > 2.5
\end{equation}

\begin{lstlisting}[style=py,caption={$z$-score detector (lines 224--229).}]
def detect_anomalies(delays, thresh=2.5):
    if len(delays) < 4:
        return []
    m = statistics.mean(delays)
    s = statistics.stdev(delays)
    return [] if s == 0 else [i for i, d in enumerate(delays) if abs((d-m)/s) > thresh]
\end{lstlisting}

\begin{figure}[H]
\centering
\begin{tikzpicture}[x=1.05cm,y=1cm,font=\scriptsize]
% axes
\draw[->,draw=slate] (0,0) -- (11.2,0) node[below left,font=\tiny]{hop index};
\draw[->,draw=slate] (0,0) -- (0,3.4) node[above right,font=\tiny,align=left]{delay (s)};
% mean and bands
\draw[draw=slate,thick] (0.2,1.2) -- (10.6,1.2);
\node[font=\tiny,text=slate,anchor=west] at (10.7,1.2) {$\mu$};
\begin{scope}[on background layer]\fill[mist] (0.2,0.75) rectangle (10.6,1.65);\end{scope}
\draw[draw=amber,dashed] (0.2,2.55) -- (10.6,2.55);
\node[font=\tiny,text=amber,anchor=west] at (10.7,2.55) {$\mu+2.5\sigma$};
% normal points
\foreach \x/\y in {0.7/1.15, 1.5/1.32, 2.3/1.05, 3.1/1.4, 3.9/1.12, 4.7/1.28,
                   6.3/1.18, 7.1/1.36, 7.9/1.08, 8.7/1.25, 9.5/1.3}
  {\fill[teal] (\x,\y) circle (2.1pt);}
% anomaly
\fill[coral] (5.5,2.95) circle (3.2pt);
\draw[draw=coral,thick] (5.5,1.2) -- (5.5,2.9);
\node[font=\tiny,text=coral,anchor=west,align=left] at (5.7,3.12)
  {$z = 4.1$ \ \textbf{FLAGGED}\\ possible delay injection};
\node[font=\tiny,text=slate,anchor=west] at (0.2,-0.55)
  {Threshold 2.5 catches roughly 99\,\% of injections at a false-positive rate below 1.2\,\% on a normal delay distribution.};
\end{tikzpicture}
\caption[Z-score delay anomaly detection]{\textbf{Detecting a delay-injection attempt.}
Ordinary variation from mobility and multipath stays within a couple of standard deviations.
An artificially delayed packet stands out as a statistical outlier and is flagged, even
though it may still be inside the $\Delta T$ window and would otherwise be accepted.}
\label{fig:zscore}
\end{figure}

%----------------------------------------------------------------------
\section{E7 --- Extended metrics and comparison charts}

The original produced three plots showing delay, energy and communication cost against node
count, with no external reference. We extended this to seven, adding:

\begin{itemize}
\item comparison bar charts against all five prior schemes from the paper's Tables~IV and~V,
      on a log scale for computational cost;
\item reference lines on the energy and communication plots so the reader can see the margin
      rather than infer it;
\item a network topology plot showing the failed node and the recovered route;
\item \textbf{throughput} (authentications per second $= 1/\bar{d}$), a metric absent from
      the base paper;
\item \textbf{battery level per node} after the simulation, showing which nodes bear the
      cost.
\end{itemize}

All seven are presented and discussed in Chapter~\ref{ch:results}.
